\documentclass[../../main.tex]{subfiles}
\begin{document}

{\bf [解]}
$C: y=\log x$，$(\log x)'=\dfrac{1}{x}$ より，$C$ 上 $x=t$ での接線は
\begin{align}
\ell_t: y=\dfrac{1}{t}(x-t)+\log t
\end{align}
だから
\begin{align}
\ell_a: y=\dfrac{1}{a}(x-a)+\log a, \qquad \ell_c: y=\dfrac{1}{c}(x-c)+\log c
\end{align}
とおり，$\ell_a,\ell_c$ の交点 $R$ の $x,y$ 座標$x_R, y_R$は
\begin{align}
\begin{cases}
 x_R=\dfrac{ac(\log c-\log a)}{c-a} \\
 y_R=\dfrac{c\log c-a\log a}{c-a}
\end{cases}\label{eq:1}
\end{align}
である．

\begin{figure}[htb]
\begin{center}
\begin{tikzpicture}[scale=2, >=stealth]
  % パラメータ設定 (a = 1.2, c = 2.6)
  \pgfmathsetmacro{\a}{1.2}
  \pgfmathsetmacro{\c}{2.6}
  
  % 交点 R の実際の計算 (a=1.2, c=2.6)
  % x_R = a*c*(ln(c)-ln(a))/(c-a) ≒ 1.723
  % y_R = (c*ln(c)-a*ln(a))/(c-a) ≒ 0.618
  \pgfmathsetmacro{\xr}{1.723}
  \pgfmathsetmacro{\yr}{0.618}

  % 1. 領域 S (曲線 y=log x と x軸, x=a, x=c で囲まれた領域) の塗りつぶし
  \begin{scope}
    \fill[pattern=north east lines, pattern color=blue!70!black]
      (\a, 0) -- plot[domain=\a:\c, samples=50] (\x, {ln(\x)}) -- (\c, 0) -- cycle;
  \end{scope}

  % 2. 領域 T (三角形 P'Q'R) の塗りつぶし
  \begin{scope}
    \fill[pattern=north west lines, pattern color=red!70!black]
      (\a, 0) -- (\xr, \yr) -- (\c, 0) -- cycle;
  \end{scope}

  % 3. 座標軸
  \draw[->] (-0.3,0) -- (4.0,0) node[right] {$x$};
  \draw[->] (0,-0.8) -- (0,2.0) node[above] {$y$};
  \node[below left] at (0,0) {$O$};

  % 4. 曲線 C: y = log x
  \draw[domain=0.3:3.6, smooth, thick] plot (\x, {ln(\x)});
  \node[above] at (3.5, {ln(3.5)}) {$C: y=\log x$};

  % 5. 接線 l_a, l_c
  \draw[thick, blue!80!black, domain=0.4:3.2] plot (\x, {(\x-\a)/\a + ln(\a)});
  \draw[thick, red!80!black, domain=0.8:3.5] plot (\x, {(\x-\c)/\c + ln(\c)});
  
  \node[blue!80!black, above left, font=\small] at (0.5, { (0.5-\a)/\a + ln(\a) }) {$\ell_a$};
  \node[red!80!black, above right, font=\small] at (3.3, { (3.3-\c)/\c + ln(\c) }) {$\ell_c$};

  % 6. 補助線 (x=a, x=c の破線)
  \draw[dashed, gray] (\a, 0) -- (\a, {ln(\a)});
  \draw[dashed, gray] (\c, 0) -- (\c, {ln(\c)});

  % 7. 点のプロットとラベル
  \fill (\a, 0) circle (1.2pt) node[below] {$P'(a,0)$};
  \fill (\c, 0) circle (1.2pt) node[below] {$Q'(c,0)$};
  \fill (\a, {ln(\a)}) circle (1.2pt) node[above left=-1pt] {$P$};
  \fill (\c, {ln(\c)}) circle (1.2pt) node[above left=-1pt] {$Q$};
  \fill (\xr, \yr) circle (1.2pt) node[above=2pt] {$R$};

  % 8. 領域ラベル (S, T)
  \node[blue!90!black, fill=white, inner sep=1pt, font=\small] at (1.9, 0.25) {$S$};
  \node[red!90!black, fill=white, inner sep=1pt, font=\small] at (1.7, 0.45) {$T$};

\end{tikzpicture}
\end{center}
\caption{面積 $S$（積分領域）と 面積 $T$（$\triangle P'Q'R$）}
\end{figure}

まず$S$の面積は
\begin{align}
 S 
  &=\int_a^c\log x\,dx \\
  &=\left[x(\log x-1)\right]_a^c \\
  &=c(\log c-1)-a(\log a-1) \quad \cdots \text{①} \label{eq:2}
\end{align}
と書ける．次に$\triangle P'Q'R$の面積$T$は$P'Q'$を底辺を見立てると
\begin{align}
 T=\dfrac{1}{2}(c-a) y_R
\end{align}
だから，\cref{eq:1}を代入して
\begin{align}
 T
  &= \dfrac{1}{2}(c-a) \dfrac{c\log c-a\log a}{c-a} \\
  &= \dfrac{c\log c-a\log a}{2} \label{eq:3}
\end{align}
である．\cref{eq:2,eq:3}より求める比は
\begin{align}
S:T=\bigl\{c(\log c-1)-a(\log a-1)\bigr\}:\dfrac{1}{2}(c\log c-a\log a)
\end{align}
である．


\end{document}
